\onecolumn

\section{Major Revision}
\begin{itemize}
    \item [Done, Jonas] ip:port (x) combinations vs. ip-port combinations vs. ip, port combinations
    \item [Done, 22.05 15:25, @Gurur] IP vs. ip
    \item [Done, 22.05 22:00, @Gurur] naiive vs. naive (x) vs. naïve
    \item [Done, 22.05 22:15, @Gurur] Internet vs. internet, not sure but used Internet
    \item [Done, 22.05 22:20, @Gurur] RFC-1337, RFC1337, RFC 1337 (x).
    \item [Done, 22.05 11:54 @Jonas - 1,2 times strong does still occur, i contacted gustavo about this, he aims to fix it, maybe some one else can have a look here, too?]strong/weak TLS vs. recommended/not recommended TLS.
    \item [Done, @Jonas] 1.234 vs. 1,234 vs. 1234
    \item [Done, @Jonas] Listing 1
    \item [Done, @Jonas] Bibliography hat Typos, alle durchgehen.
    \item [Done, @Gurur] The numbers of sections 4.1 (scanning) and 4.2 (confirmation) are not matching.
    \item merge, shorten, and fix numbers in 5.1 and 5.2 (maybe later work, because the reviewers didn't request any changes about this section)
    \item [Done, @Jonas] The phrasing in section 4.3 is hard to understand (I could not get the meaning to be able to rewrite it).
    \item [Done, @Jonas] layout of tables 14 til 17. Maybe merge them all into one big table.
    \item [Done, 22.05 15:15, @Gurur] root\_dse, Root DSE, root DSE (WOULD SUGGEST THIS), root dse, ...
    \item [Done, @Gurur] MitM vs. PitM
    \item [Done, @Gurur] \eg vs e.g.
\end{itemize}

\textbf{Before submission}
\todo{All paper submissions, including Major Revisions, should be at 
most 13 typeset pages, excluding bibliography and well-marked 
appendices.}
\subsection*{Typos and Stuff}
\begin{itemize}
    \item \todo{IMPORTANT:} Purge any author identifying information! No Gurur, Jonas, Gustavos, FHMS, in the text and Figures! This will lead to auto-rejects from USENIX!!
    \item STARTTLS vs StartTLS vs Start TLS - Correct one is StartTLS for LDAP (and STARTTLS for SMTP, POP3, and IMAP)~\rh{I defined a newcommand for that in acronyms.tex}
    \item Consistent use of acronyms
    \item Bibliography is missing numerous attributes
    \item Clarify early on that we do not look at LDAP DDoS amplification (UDP).
    \item TCP, UDP lower or upper case? \rh{always upper}
    \item Table: Reasons for invalid X.509 certificate chain -> PII to personal data (if image is used)
    \item Decided Names: paper title, framework - and their correct spelling
    \item Unify numbering in tables and text (10000 vs. 10,000)
%    \item unauthenticated, vs simple bind: two different things, what do you mean?
    \item Landscape Pipeline: Analyse => Analyze
%    \item remove Gustavo's name in the reference (code ref)
    \item remove refs on repos of authors? 
    \item naming contexts not explained anywhere
    \item e.g., and et al. (macros) italic?
    \item dots in paragraph headings?
    \item table headers bold vs plain
%    \item Remove Table 13: Grouping of LDAP attributes containing personal data?
    \item Remove Table 16: Recommended and non-recommended cipher suites on LDAP hosts supporting TLS?
\end{itemize}

\textbf{Checked 08.02. - 21.00}
\begin{itemize}
    \item Landscape Pipeline: rename Validation Module to Confirmation Module

    \item clear text (dont) vs. clear-text (dont), vs cleartext (do)
    \item Analyse vs. Analyze => US: Analyze
    \item LDAP v3 vs LDAPv3 (same with v2) => LDAPv2
    \item ZMap (not ZMAP) => fixed
    \item use Person in the Middle instead of Man in the Middle or better yet: Meddler in the Middle to keep the acronym
    \item Spelling use US-style: Analyze, color, behavior
    \item plain text vs plaintext => plaintext in cryptocontext. Plain text in ASCII/Text context
    \item Cites with multiple items (do) vs. multiple cites (don't)
    \item Cites need non-breaking spaces before them($\sim$) %asdf~\cite{ref}
    \item ref vs. cref vs. autoref  => CREF
    \item consistently use \eg and \etal as macros
    \item Optimize table 12 'Grouping of LDAP attributes containing personal data' wrt white spaces
\end{itemize}


\vfill
\pagebreak
\twocolumn